\section{做题}
\begin{enumerate}
  \item \zx{主体错、行为错、应当/可以错}。
  \item 逻辑词。
        \begin{enumerate}
          \item 可以 = 能够 \bdy 可能。
          \item 选项出现“如 xx”，不要质疑假设能否成立。
          \item \zd{程序法里的应当 = 公检法没有自由裁量权}。
          \item 不得 = 绝无可能、绝对不。
          \item 应当 = 必须。
          \item 可促成、可积极促成。\zx{注意加限定修饰词的，不是说都是错，而是要谨慎再谨慎}！
        \end{enumerate}
  \item 一定要梗概。\al{在审判阶段，\xhx{值班律师}可以出庭进行\xhx{辩论} \cuo}。
  \item 区分不同主体，不同阶段。
        \begin{enumerate}
          \item \al{县公安局以甲涉嫌过失致人死亡\xhx{立案侦查}，后随即\xhx{撤案}，乙的近亲属的诉讼权利。\scx{就撤销向县公安局复议}。只能公转自了}。
          \item 注意执行阶段，已无嫌疑人/被告人身份。
          \item 不要忘了最高法、最高检。
          \item 法院除了审判长、还有\zx{院长}。
        \end{enumerate}
  \item 以偏概全。\al{外籍甲只能委托具有中国律师资格并取得执业证书的律师作为辩护人}。做梗概。
        \begin{enumerate}
          \item \al{只能委托律师作为辩护人}。
          \item \al{仅仅、都要、一律、必须}。
        \end{enumerate}
  \item 选项看起来冠冕堂皇很有理由，应当就是应当，没有例外，一般应当才有例外。
        \begin{enumerate}
          \item \al{速裁程序，某县法院经过审理认为案件复杂，可以在庭后听取控辩双方意见后\xhx{定期宣判}}。
        \end{enumerate}
  \item 选项单独拿出来看没问题，结合题干肯定错。\zx{对这类表述一定一定要多看几次题目}！！
  \item 正话反说、张冠李戴、答非所问/无关选项、声东击西、无中生有、故弄玄虚。
\end{enumerate}
